% Nouveau EGA V (cote 25): archive cote 25, site page 150 (PDF page 151).
% typescript. Draft compared: none.
% AI-transcribed by Claude Opus 5.5 (claude-opus-5-5, adaptive thinking, effort high), Message Batches API, batch ega25b (msgbatch_01A7n5YuXubUPNK1zsepi8RG), 2026-11-13; not yet checked by eye.
\scanpage{25/150}
\mspage{\hand{11}}
\placedmargin{\unsure{$\diagup$}}
% NOTE: A diagonal handwritten stroke runs from the top edge above conditions to the left margin beside k >=1. It stops before Proposition 9. Its cancellation or working-mark function remains uncertain; the text is retained with an uncertain marginal diagonal sign.
d'ailleurs de vérifier ces dernières conditions pour les fibres réduites
en les différents idéaux premiers (ou même seulement maximaux) de A/I .
\uline{Proposition} 8. Pour que l'augmentation A$\hand{\rightarrow}$A/I soit de type quadratique
ordinaire, il f et s que les conditions suivantes soient vérifiées:
1º)  Les $gr^{n}_{I}(A)$ sont projectifsde type fini sur A/I .

2º) Si d désigne le rang de $gr^{1}_{I}(A) = I/I^{2}$ en un $\underline{p} \hand{\in}$Spec(A/I),
alors le rang de $gr^{n}_{I}(A)$ en ce même point est égal (ou simplement
$\hand{\geqslant}$ ) à $\binom{n+d}{d} - \binom{n-2+d}{d}$ .

3º) (1º et 2º impliquent déjà que $\underline{K}$ = Ker $\underline{\theta}^{\overset{\hand{\unsure{6|\sigma}}}{2}}$ est un Module
inversible localement facteur direct dans  Sym$^{2}$(I/I$^{2}$) ). Le sous
Module $\underline{K}$ = Ker $\theta^{2}$ de Sym$^{2}$(I/I$^{2}$) est non dégénéré.

La démonstration est immédiate à partir du lemme suivant; 

 \uline{Lemme} 8.1. Soient B un anneau, M un B-module projectif de
type fini, $\hand{\phi} \hand{\in}$ Sym$^{2}$(M) un élément ``non dégénéré'' , alors la multiplica-%
tion par $\hand{\phi}$
\[ \mathrm{Sym}^{\overset{\hand{\illegible{1}}}{n\hand{\unsure{-2}}}}(M) \hand{\longrightarrow} \mathrm{Sym}^{n}(M) \]
est un isomorphisme sur un sous-module facteur direct, donc le conoyau
est un Module projectif de rang $\binom{n+d}{d} - \binom{n-2+d}{d}$ si M est de rang d.

On est immédiatement ramené au cas où B est le spectre d'un corps,
où c'est trivial (en fait, le ``non dégénéré'' nous a servi uniquement
\placedmargin{||} pour assurer que les $\hand{\phi}(x)$, $x \hand{\in}$ Spec(B), sont tous $\hand{\neq} 0$, et on pourrait
formuler le lemme plus raisonnablement en remplaçant 2 par un entier
k $\hand{\geqslant}$1 quelconque, sans parler de non dégénérescence).

\uline{Proposition} 9. Soit A un anneau augmenté vers A/I, à augmentation
régulière, I/I$^{2}$ étant projectif de rang n sur A/I . Soit E un A-module

\mspage{\hand{150}}
